% Test-case specification TC-0003 (CTest test TC-0003-helpers-clock).
\documentclass[letterpaper]{article}

\usepackage[letterpaper,margin=1in]{geometry}
\emergencystretch=3em

% Compile from the repository root so this path resolves to
% the pinned WSP submodule.
\input{wsp/testing/test-case-library.tex}

\begin{document}

\def\TCID{TC-0003}
\def\TCTitle{%
  Date\$, Time\$ and the present date and time}
\def\TCPurpose{%
  Verify that \texttt{rtcGetDateBstr}, \texttt{rtcGetTimeBstr},
  \texttt{rtcGetDateVar}, \texttt{rtcGetTimeVar} and
  \texttt{rtcGetPresentDate} return the current date and time in the
  reference formats.}
\def\TCPriority{\TCPriorityHigh}

\def\TCRequirementRef{%
  REQ-0004}
\def\TCDesignRef{%
  Not applicable.}
\def\TCFeatureRef{%
  \texttt{tests/\allowbreak{}vbrun\_\allowbreak{}helper\_\allowbreak{}tests.\allowbreak{}cpp}
  (family \texttt{clock}),
  \texttt{src/\allowbreak{}helpers\_\allowbreak{}clock.\allowbreak{}cpp}}
\def\TCTestTechnique{%
  Requirements-based testing against the system clock.}

\def\TCPreconditions{%
  The library and the test program are built
  (\texttt{scripts/\allowbreak{}Invoke-Build.\allowbreak{}ps1}) and the test
  is registered with CTest.}
\def\TCEnvironment{%
  Windows with the toolchain of
  \texttt{scripts/\allowbreak{}Invoke-Build.\allowbreak{}ps1}; the local
  gate builds and runs x64 Debug, CI builds and runs the other
  architectures. The library is built from this repository in
  \texttt{out/\allowbreak{}build}.}
\def\TCAssumptions{%
  Not applicable.}

\def\TCInputData{%
  The system clock, read before and after each call.}
\def\TCInitialState{%
  No state from earlier runs is needed; the test makes what it needs.}

\def\TCProcedure{%
  \begin{enumerate}
    \item CTest runs \texttt{vbrun\_\allowbreak{}helper\_\allowbreak{}tests
      \textless{}wfcvbrun.\allowbreak{}dll\textgreater{} clock}.
    \item Three times: read the clock, call \texttt{rtcGetDateBstr} and
      \texttt{rtcGetTimeBstr}, read the clock again.
    \item Call \texttt{rtcGetDateVar}, \texttt{rtcGetTimeVar} and
      \texttt{rtcGetPresentDate} between two clock readings.
    \item CTest records the exit code and, on failure, the program's output
      (\texttt{--output-on-failure}).
  \end{enumerate}}
\def\TCExpectedResult{%
  The text of \texttt{Date\$} is \texttt{mm-dd-yyyy} and the text of
  \texttt{Time\$} is \texttt{hh:mm:ss} of the clock reading before or after
  the call. Each Variant is of type \texttt{VT\_\allowbreak{}DATE} and its
  serial equals the date, the time or both of one of the two readings. The
  program exits 0.}
\def\TCPostConditions{%
  No file or registry state remains that later tests depend on.}

\TCTemplate

\end{document}
